\begin{table}[!htbp]
    \centering
    \caption{Alasan pemilihan dan kelemahan yang diakui pada keputusan teknis.}
    \label{tab:tab3.7_alasan_keputusan_teknis}
    \renewcommand{\arraystretch}{1.10}
    \begin{tabular}{@{}p{2.8cm}p{3.2cm}p{5.0cm}p{4.0cm}@{}}
        \hline
        \textbf{Keputusan} & \textbf{Alternatif terpilih} & \textbf{Alasan pemilihan} & \textbf{Kelemahan yang diakui} \\
        \hline
        Model prediksi & Gradient Boosting & Dipilih berdasarkan ketepatan kelas kondisi dan keterterapan pada lingkungan komputasi penelitian. Keluaran dapat digunakan untuk jalur regresi dan klasifikasi dalam keluarga model yang sama. & Kepekaan terhadap hipoglikemia pada horizon 60 menit lebih rendah daripada Random Forest pada evaluasi yang digunakan untuk pengambilan keputusan. \\
        Model bahasa & Layanan awan & Tidak membebani perangkat lokal dan tersedia pada skema penggunaan bertingkat yang sesuai dengan konteks penelitian. & Memerlukan konektivitas jaringan. \\
        Model embedding & Model lokal & Dapat berjalan tanpa akselerator grafis dan tidak menimbulkan biaya per permintaan. & Kemampuan penelusuran lintas bahasa bergantung pada karakteristik model embedding yang digunakan. \\
        Pengambilan pengetahuan & RAG atas korpus tidak terstruktur & Tidak menuntut pembangunan ontologi maupun integrasi rekam medis sehingga lebih sesuai dengan batasan infrastruktur. & Mutu hasil bergantung pada cakupan dan kualitas korpus. \\
        Pengklasifikasi kondisi & Pengklasifikasi terpisah & Kelas kondisi dapat dioptimalkan langsung untuk kebutuhan pembentukan kueri dan tidak hanya bergantung pada ambang keluaran regresi. & Menambah satu model yang perlu dilatih dan dipelihara. \\
        Orkestrasi alur & Penulisan sebagai kode program & Mendukung pemisahan logika domain dan dapat berjalan tanpa server khusus. & Memerlukan penulisan dan pemeliharaan kode lebih banyak dibandingkan perkakas visual. \\
        \hline
    \end{tabular}
\end{table}
